\subsection{Caso de Estudio 1.2: Tetris Web}
\label{subsec:tetris_init}

La inicialización de Tetris se ejecutó cinco veces de forma independiente, una por herramienta, con una misma semilla en inglés (versión 1.1; SHA-256 cuyo prefijo es \texttt{2AE084C9C64F}). Se eligió la opción A de la matriz experimental: tablero de $10\times20$, siete tetrominós en orientación inicial, gravedad de 500 ms, desplazamientos izquierda/derecha/abajo, colisiones, fijación, eliminación de líneas, fin de partida y reinicio. El servidor FastAPI conserva el estado autoritativo en memoria y React con TypeScript representa ese estado. Rotación, puntuación, persistencia, clasificación y demás ampliaciones quedaron fuera de esta etapa. Las cinco implementaciones se conservaron en directorios separados y fueron verificadas mediante compilación del cliente, pruebas disponibles y análisis de SonarQube. Para BMAD se reporta únicamente la ejecución con planificación completa: \textit{product brief}, PRD, arquitectura, épicas e historias, construcción y corrección de alcance. Se utilizó BMAD Method 6.12.0 con los módulos \texttt{core} y \texttt{bmm} y la integración Gemini; la semilla y el alcance permanecieron iguales a los de las otras cuatro herramientas.

La Tabla \ref{tab:versiones_tetris_init} precisa las versiones utilizadas en estas ejecuciones, que difieren de algunas versiones de referencia de la Tabla \ref{tab:stack_tecnologico}. En particular, la ejecución de Tessl utilizó la CLI vigente como gestor de contexto, con plugins públicos para FastAPI y React y su integración MCP con Gemini; no instaló el \textit{tile} SDD 2.0.1 presentado como referencia teórica en el capítulo. Esta diferencia metodológica debe conservarse al interpretar sus resultados.

\begin{table}[H]
\centering
\caption{Versiones efectivas de las herramientas en Tetris Etapa 1}
\label{tab:versiones_tetris_init}
\begin{tabularx}{\textwidth}{l Y}
\toprule
\textbf{Componente} & \textbf{Versión observada} \\
\midrule
OpenSpec & 1.13.2 \\
GitHub Spec Kit & 1.0.12 \\
Tessl CLI & 0.108.0; plugins de contexto públicos y MCP \\
BMAD Method & 6.12.0 \\
Kiro CLI & 2.25.0 (modo Auto) \\
Gemini CLI & 0.60.0 (OpenSpec, Spec Kit, Tessl y BMAD) \\
SonarQube & Community Build 26.9.0.129388 \\
\bottomrule
\end{tabularx}
\end{table}

\subsubsection{Dimensión A: Esfuerzo Computacional y Costos}

La Tabla \ref{tab:consumo_tetris_init_costos} sigue las categorías empleadas para la calculadora. Los cuatro flujos basados en Gemini CLI utilizaron Vertex AI; Kiro CLI se ejecutó en modo \textit{Auto}. Los valores precedidos por $\geq$ son el tráfico mínimo recuperable de los registros por ejecución, no el consumo facturado ni un conteo completo de la etapa: falta el registro de tokens de una ejecución de OpenSpec, una de Spec Kit y tres de Tessl. Se distingue la entrada total de la porción cacheada para no cobrar dos veces los mismos tokens. La estimación económica aplica la fórmula y las tarifas indicadas antes en este capítulo únicamente al tráfico observado; no equivale a una factura de GCP, que puede estar cubierta por créditos del proyecto. El protocolo ACP conservado de Kiro no registra tokens ni créditos consumidos de forma verificable, por lo que tampoco permite calcular su costo.

\begin{table}[H]
\centering
\caption{Tráfico observable y costo teórico del modelo -- Inicialización (Tetris Web)}
\label{tab:consumo_tetris_init_costos}
\small
\begin{tabularx}{\textwidth}{l Y Y Y Y Y}
\toprule
\textbf{Herramienta} & \textbf{Input tokens} & \textbf{Input cacheado} & \textbf{Output tokens} & \textbf{Créditos} & \textbf{Costo observable (USD)} \\
\midrule
OpenSpec & $\geq$ 5.816.534 & $\geq$ 4.217.424 & $\geq$ 39.707 & N/A & $\geq 3,3886$ \\
Spec Kit & $\geq$ 14.171.471 & $\geq$ 10.956.018 & $\geq$ 90.198 & N/A & $\geq 7,2784$ \\
Tessl & $\geq$ 5.952.918 & $\geq$ 3.896.072 & $\geq$ 41.324 & N/A & $\geq 4,0416$ \\
BMAD & 17.653.626 & 12.533.513 & 237.579 & N/A & $11,6984$ \\
Kiro & N/D & N/D & N/D & N/D & N/D \\
\bottomrule
\multicolumn{6}{p{\dimexpr\textwidth-2\tabcolsep}}{\footnotesize \textit{Nota:} N/A significa que la plataforma no usa esa unidad; N/D indica que la ejecución no dejó un dato verificable. En OpenSpec, Spec Kit y Tessl faltan 1/6, 1/6 y 3/8 registros de tokens, respectivamente. Los mínimos se refieren a sesiones registradas y no a la factura total.}
\end{tabularx}
\end{table}

\begin{figure}[H]
\centering
\begin{tikzpicture}
\begin{axis}[
    ybar stacked, width=0.95\textwidth, height=6.3cm,
    symbolic x coords={OpenSpec,SpecKit,Tessl,BMAD}, xtick=data,
    ymin=0, bar width=17pt, enlarge x limits=0.15,
    ylabel={Costo teórico observado (USD)},
    legend style={at={(0.5,-0.2)},anchor=north,legend columns=3},
    nodes near coords={},
]
\addplot[fill=vdarkblue] coordinates {(OpenSpec,2.3987) (SpecKit,4.8232) (Tessl,3.0853) (BMAD,7.6802)};
\addplot[fill=darkblue] coordinates {(OpenSpec,0.6326) (SpecKit,1.6434) (Tessl,0.5844) (BMAD,1.8800)};
\addplot[fill=lightblue] coordinates {(OpenSpec,0.3574) (SpecKit,0.8118) (Tessl,0.3719) (BMAD,2.1382)};
\legend{Entrada nueva,Entrada cacheada,Salida}
\end{axis}
\end{tikzpicture}
\caption{Desglose del costo teórico del tráfico de tokens recuperable. Las barras incompletas de OpenSpec, Spec Kit y Tessl representan mínimos; Kiro carece de datos verificables.}
\label{fig:costos_tetris_init}
\end{figure}

Entre las sesiones medibles, BMAD registra la mayor entrada observada ($17,65$ millones de tokens) y el mayor costo teórico del tráfico ($11,70$ USD). Esta cifra suma entrada nueva, entrada cacheada y salida con las tarifas declaradas; no corresponde a una factura de GCP. Los importes de OpenSpec, Spec Kit y Tessl son mínimos por sus registros incompletos, de modo que la comparación económica no identifica un costo final ordenado para las cinco herramientas. Kiro permanece fuera de esta comparación monetaria por falta de créditos trazables, aunque su ejecución y sus artefactos son observables.

\begin{table}[H]
\centering
\caption{Esfuerzo operativo registrado -- Inicialización (Tetris Web)}
\label{tab:consumo_tetris_init_esfuerzo}
\begin{tabularx}{\textwidth}{l Y Y Y Y}
\toprule
\textbf{Herramienta} & \textbf{Solicitudes registradas} & \textbf{Tiempo de agente (s)} & \textbf{Llamados a herramientas} & \textbf{Sesiones sin tokens} \\
\midrule
OpenSpec & 6 & 3.065 & $\geq$ 85 & 1 \\
Spec Kit & 6 & 3.359 & $\geq$ 162 & 1 \\
Tessl & 8 & 4.901 & $\geq$ 86 & 3 \\
BMAD & 9 & 3.299,511 & $\geq$ 342 & 2 \\
Kiro & 1 & 546 & 57 & N/A \\
\bottomrule
\multicolumn{5}{p{\dimexpr\textwidth-2\tabcolsep}}{\footnotesize \textit{Nota:} Solicitud = invocación registrada del agente, no cada mensaje interno del modelo. El tiempo suma la duración de esas invocaciones, incluso intentos fallidos; excluye instalaciones, pruebas independientes y escaneo. Los llamados mínimos de Gemini solo cubren salidas JSON con estadística de herramientas; en Kiro se contaron eventos \texttt{tool\_call} del protocolo ACP. Estas definiciones difieren de las interacciones por mensaje usadas para Calculadora, por lo cual sus valores absolutos no deben compararse directamente entre casos.}
\end{tabularx}
\end{table}

\begin{figure}[H]
\centering
\begin{tikzpicture}
\begin{axis}[
    ybar, width=0.9\textwidth, height=5.7cm,
    symbolic x coords={OpenSpec,SpecKit,Tessl,BMAD,Kiro}, xtick=data,
    ymin=0, bar width=17pt, enlarge x limits=0.13,
    ylabel={Tiempo registrado (minutos)},
    nodes near coords, every node near coord/.append style={font=\scriptsize},
]
\addplot[fill=darkblue] coordinates {(OpenSpec,51.08) (SpecKit,55.98) (Tessl,81.68) (BMAD,54.99) (Kiro,9.10)};
\end{axis}
\end{tikzpicture}
\caption{Duración acumulada de las invocaciones del agente en Tetris, sin tiempo de instalación ni validación externa.}
\label{fig:tiempo_tetris_init}
\end{figure}

Tessl exigió ocho solicitudes y $81,7$ minutos de agente, el máximo del grupo; su ejecución incluyó respuestas de Vertex con error 429 o planes sin archivos, y después una corrección para restringir los controles a las flechas exigidas por la semilla. Kiro terminó su flujo nativo en una invocación de $9,1$ minutos, aunque ese tiempo no representa el esfuerzo total de preparación del entorno. OpenSpec y Spec Kit requirieron seis invocaciones cada uno. En Spec Kit se corrigió una prueba de aceptación imposible para una pieza I vertical sin rotación; en BMAD se preservaron 30 requisitos funcionales y se eliminaron metas de capacidad y latencia añadidas por el agente sin estar en la semilla. Sus nueve invocaciones incluyen un arranque fallido por permisos y una reanudación vacía sin tokens, además de las fases de planificación y construcción. Por tanto, la duración refleja tanto la herramienta como la necesidad de corrección y las condiciones de servicio; no debe interpretarse como una medición aislada de velocidad de generación.

\subsubsection{Dimensión B: Calidad del Código Generado (SonarQube)}

Las cinco variantes superaron su \textit{quality gate} en SonarQube Community Build 26.9.0.129388. La Tabla \ref{tab:sonar_tetris_init} usa \texttt{ncloc} para LOC, las categorías clásicas \texttt{bugs}, \texttt{vulnerabilities} y \texttt{code\_smells}, y los ratings Security/Reliability/Maintainability obtenidos de la API. Un valor 1 equivale a A y 3 a C. El \textit{quality gate} aprobado no implica ausencia de hallazgos en el código total: SonarQube puede evaluar condiciones de código nuevo distintas de las cifras globales aquí presentadas.

\begin{table}[H]
\centering
\caption{Métricas estáticas SonarQube -- Inicialización (Tetris Web)}
\label{tab:sonar_tetris_init}
\scriptsize
\begin{tabularx}{\textwidth}{l Y Y Y Y Y Y Y Y}
\toprule
\textbf{Herramienta} & \textbf{LOC} & \textbf{Bugs} & \textbf{Vulns.} & \textbf{Smells} & \textbf{Deuda (m)} & \textbf{Comp. cog.} & \textbf{Dupl.} & \textbf{Rating (S/R/M)} \\
\midrule
OpenSpec & 898 & 0 & 1 & 15 & 81 & 76 & 0\% & C/A/A \\
Spec Kit & 748 & 0 & 2 & 7 & 35 & 78 & 0\% & C/A/A \\
Tessl & 788 & 0 & 1 & 13 & 71 & 69 & 0\% & C/A/A \\
BMAD & 669 & 0 & 1 & 13 & 84 & 97 & 0\% & C/A/A \\
Kiro & 497 & 0 & 0 & 10 & 38 & 45 & 0\% & A/A/A \\
\bottomrule
\end{tabularx}
\end{table}

Para construir un radar comparable en estructura con la Figura \ref{fig:radar_calc_init}, se usaron los mismos cinco ejes y se normalizaron por separado los resultados de Tetris. En cada eje donde los valores difieren, una magnitud menor recibe mejor puntuación mediante $s(x)=2+8(x_{\max}-x)/(x_{\max}-x_{\min})$; si todos los valores son iguales, se asigna $s=10$ a las cinco variantes. La escala es relativa a este caso de estudio y no constituye una calificación SonarQube ni debe compararse numéricamente con la escala normalizada para Calculadora.

\begin{figure}[H]
\centering
\begin{tikzpicture}[font=\footnotesize]
  \foreach \r in {0.6,1.2,1.8,2.4,3.0} {
    \draw[gray!30] (0:\r) -- (72:\r) -- (144:\r) -- (216:\r) -- (288:\r) -- cycle;
  }
  \foreach \ang in {0,72,144,216,288} {\draw[gray!55] (0,0) -- (\ang:3.0);}
  \node[anchor=west] at (0:3.35) {Bugs};
  \node[anchor=south] at (72:3.35) {Vulnerabilidades};
  \node[anchor=east] at (144:3.35) {Code Smells};
  \node[anchor=east,align=right] at (216:3.35) {Complejidad\\cognitiva};
  \node[anchor=north] at (288:3.35) {Duplicación};

  \draw[line width=1.3pt,color=vdarkblue] (0:3) -- (72:1.8) -- (144:1.2) -- (216:0.6) -- (288:3) -- cycle; % BMAD
  \draw[line width=1.3pt,color=darkblue] (0:3) -- (72:1.8) -- (144:0.6) -- (216:1.569) -- (288:3) -- cycle; % OpenSpec
  \draw[line width=1.3pt,color=blue] (0:3) -- (72:0.6) -- (144:3) -- (216:1.477) -- (288:3) -- cycle; % Spec Kit
  \draw[line width=1.3pt,color=purple] (0:3) -- (72:3) -- (144:2.1) -- (216:3) -- (288:3) -- cycle; % Kiro
  \draw[line width=1.3pt,color=lightblue] (0:3) -- (72:1.8) -- (144:1.2) -- (216:1.892) -- (288:3) -- cycle; % Tessl
\end{tikzpicture}
\par\smallskip
{\footnotesize
\textcolor{vdarkblue}{\rule[0.5ex]{0.35cm}{1.3pt}} BMAD\quad
\textcolor{darkblue}{\rule[0.5ex]{0.35cm}{1.3pt}} OpenSpec\quad
\textcolor{blue}{\rule[0.5ex]{0.35cm}{1.3pt}} Spec Kit\quad
\textcolor{purple}{\rule[0.5ex]{0.35cm}{1.3pt}} Kiro\quad
\textcolor{lightblue}{\rule[0.5ex]{0.35cm}{1.3pt}} Tessl}
\caption{Radar de calidad estática normalizada -- Inicialización Tetris. El exterior representa menos hallazgos o menor complejidad dentro de este caso. Bugs y duplicación valen 10 para todos porque sus valores observados fueron cero en las cinco variantes.}
\label{fig:radar_tetris_init}
\end{figure}

\begin{table}[H]
\centering
\caption{Densidad de hallazgos y cobertura observada -- Inicialización (Tetris Web)}
\label{tab:sonar_tetris_init_adicional}
\begin{tabularx}{\textwidth}{l Y Y Y Y}
\toprule
\textbf{Herramienta} & \textbf{Hallazgos clásicos} & \textbf{Issues/KLOC} & \textbf{Cobertura global} & \textbf{Pruebas backend exitosas} \\
\midrule
OpenSpec & 16 & 17,82 & 60,2\% & 13/13 \\
Spec Kit & 9 & 12,03 & 62,0\% & 18/18 \\
Tessl & 14 & 17,77 & 60,4\% & 9/9 \\
BMAD & 14 & 20,93 & 63,2\% & 13/13 \\
Kiro & 10 & 20,12 & 72,6\% & 39/39 \\
\bottomrule
\multicolumn{5}{p{\dimexpr\textwidth-2\tabcolsep}}{\footnotesize \textit{Nota:} Hallazgos clásicos $=$ bugs $+$ vulnerabilidades $+$ \textit{code smells}; Issues/KLOC $=$ hallazgos clásicos $/\,(\mathrm{ncloc}/1000)$. La cobertura se informa como evidencia complementaria, sin utilizarla para puntuar esta etapa según el protocolo de las Etapas 1 y 2. La cantidad de pruebas varía entre variantes y no constituye por sí sola una tasa de aceptación funcional común.}
\end{tabularx}
\end{table}

\begin{figure}[H]
\centering
\begin{tikzpicture}
\begin{axis}[
    ybar, width=0.9\textwidth, height=6cm,
    symbolic x coords={OpenSpec,SpecKit,Tessl,BMAD,Kiro}, xtick=data,
    ymin=0, bar width=12pt, enlarge x limits=0.15,
    ylabel={Número de hallazgos},
    legend style={at={(0.5,-0.18)},anchor=north,legend columns=2},
    nodes near coords, every node near coord/.append style={font=\scriptsize},
]
\addplot[fill=vdarkblue] coordinates {(OpenSpec,15) (SpecKit,7) (Tessl,13) (BMAD,13) (Kiro,10)};
\addplot[fill=purple] coordinates {(OpenSpec,1) (SpecKit,2) (Tessl,1) (BMAD,1) (Kiro,0)};
\legend{Code smells,Vulnerabilidades}
\end{axis}
\end{tikzpicture}
\caption{Hallazgos clásicos de SonarQube en Tetris. Ninguna variante registró \texttt{bugs} clásicos; los hallazgos se concentran en mantenibilidad y seguridad.}
\label{fig:hallazgos_tetris_init}
\end{figure}

La Figura \ref{fig:radar_tetris_init} muestra a Kiro en el extremo favorable para vulnerabilidades y complejidad cognitiva, mientras Spec Kit encabeza el eje de \textit{code smells}. Kiro produjo el menor volumen de código medido por SonarQube (497 LOC), la menor complejidad cognitiva (45) y ninguna vulnerabilidad; es la única variante con rating A/A/A. Sin embargo, sus diez \textit{code smells} repartidos en una base menor elevan su densidad a 20,12 hallazgos/KLOC, una de las más altas del conjunto. Por eso el radar de recuentos absolutos debe leerse junto con la densidad de la Tabla \ref{tab:sonar_tetris_init_adicional}. Spec Kit tiene el menor número absoluto de hallazgos (nueve) y la menor deuda técnica (35 minutos). BMAD presenta 14 hallazgos clásicos, 20,93 por mil LOC, la mayor complejidad cognitiva (97) y la mayor deuda estimada (84 minutos), aunque su \textit{quality gate} también fue aprobado. OpenSpec registra 15 \textit{code smells}, el máximo del grupo. La duplicación medida fue 0\% en las cinco variantes, por lo que no discrimina entre ellas en esta etapa. Las vulnerabilidades detectadas incluyen el uso de generación pseudoaleatoria para piezas; requieren interpretar el contexto del juego antes de equipararlas a un riesgo explotable. La cobertura global, aunque disponible, se mantiene fuera de la valoración formal de Etapa 1 conforme a la metodología del capítulo.

\subsubsection{Dimensión C: Métricas de Especificación y Documentación}

Se contabilizaron los documentos de especificación nativos de cada herramienta y sus líneas físicas; se excluyeron \textit{README}, dependencias, salidas de pruebas y evidencia dentro de \texttt{docs/}. Para concordar con la tabla de Calculadora, el ratio código/especificación divide las LOC \texttt{ncloc} de SonarQube entre las líneas físicas de especificación. También se indica el número de requisitos atómicos identificados en cada conjunto de documentos. El número de líneas mide volumen, no calidad semántica ni completitud por sí solo.

\begin{table}[H]
\centering
\caption{Artefactos de especificación -- Inicialización (Tetris Web)}
\label{tab:specs_tetris_init}
\begin{tabularx}{\textwidth}{l Y Y Y Y Y}
\toprule
\textbf{Herramienta} & \textbf{Archivos} & \textbf{Líneas} & \textbf{Densidad (lín./arch.)} & \textbf{Código/spec} & \textbf{Requisitos} \\
\midrule
OpenSpec & 4 & 442 & 110,50 & 2,03 & 30 \\
Spec Kit & 8 & 1.002 & 125,25 & 0,75 & 38 \\
Tessl & 2 & 136 & 68,00 & 5,79 & 34 \\
BMAD & 6 & 1.120 & 186,67 & 0,60 & 30 \\
Kiro & 3 & 608 & 202,67 & 0,82 & 35 \\
\bottomrule
\end{tabularx}
\end{table}

\begin{figure}[H]
\centering
\begin{tikzpicture}
\begin{axis}[
    width=0.85\textwidth, height=7cm,
    xmin=450, xmax=950, ymin=0, ymax=1250,
    xlabel={Líneas de código SonarQube (ncloc)},
    ylabel={Líneas de especificación},
    legend style={at={(0.5,-0.2)},anchor=north,legend columns=3},
]
\addplot[only marks,mark=square*,mark size=3pt,color=vdarkblue] coordinates {(898,442)};
\addplot[only marks,mark=triangle*,mark size=3pt,color=darkblue] coordinates {(748,1002)};
\addplot[only marks,mark=*,mark size=3pt,color=blue] coordinates {(788,136)};
\addplot[only marks,mark=diamond*,mark size=3pt,color=lightblue] coordinates {(669,1120)};
\addplot[only marks,mark=pentagon*,mark size=3pt,color=purple] coordinates {(497,608)};
\addplot[thick,dashed,gray,domain=497:898,samples=2] {1279.359-0.857999*x};
\legend{OpenSpec,Spec Kit,Tessl,BMAD,Kiro,Tendencia descriptiva}
\end{axis}
\end{tikzpicture}
\caption{Relación entre LOC y líneas de especificación para las cinco variantes de Tetris. La línea punteada es el ajuste lineal de mínimos cuadrados de estos cinco puntos; no permite inferir causalidad ni predecir otras etapas.}
\label{fig:scatter_tetris_init}
\end{figure}

BMAD es el caso más documentado: seis archivos principales y 1.120 líneas, con 30 requisitos funcionales trazables. El conteo excluye dos memorias de proceso y el estado YAML del sprint. Spec Kit conserva ocho archivos y 1.002 líneas; en ambas herramientas el ratio código/especificación es inferior a uno. Tessl concentra solo 136 líneas en dos archivos y alcanza el ratio mayor (5,79), pero aun así cubre 34 enunciados de requisito; su menor volumen no prueba automáticamente una mejor abstracción. Kiro concentra 608 líneas en tres archivos y la mayor densidad por archivo (202,67). El ajuste descriptivo de la Figura \ref{fig:scatter_tetris_init} desciende en esta muestra; con cinco observaciones y flujos documentales distintos, no permite inferir una relación causal. Las correcciones registradas muestran que contar requisitos no garantiza por sí solo fidelidad al alcance.

\subsubsection{Dimensión D: Validación Funcional (Provisional)}

El protocolo definitivo de aceptación funcional aún está en revisión. Como evidencia preliminar se presentan las suites automatizadas ejecutadas al cerrar la Etapa 1; las cinco compilaciones del cliente también pasaron. Spec Kit añadió cinco pruebas de \textit{frontend} y Tessl dos; las otras tres variantes no dejaron una suite de \textit{frontend} equivalente. Cada suite tiene distinto tamaño y cobertura de escenarios, de modo que el 100\% de pruebas exitosas dentro de ella no debe llamarse \textit{Functional Acceptance Rate} comparable. La tasa de aceptación, los defectos funcionales y los defectos críticos se mantienen como no determinados hasta ejecutar un mismo catálogo de casos contra los cinco productos.

\begin{table}[H]
\centering
\caption{Evidencia preliminar de validación -- Inicialización (Tetris Web)}
\label{tab:validacion_tetris_init}
\begin{tabularx}{\textwidth}{l Y Y Y Y Y}
\toprule
\textbf{Herramienta} & \textbf{Casos ejecutados} & \textbf{Exitosos} & \textbf{Fallidos} & \textbf{FAR común} & \textbf{Defectos/Críticos} \\
\midrule
OpenSpec & 13 & 13 & 0 & N/D & N/D \\
Spec Kit & 23 & 23 & 0 & N/D & N/D \\
Tessl & 11 & 11 & 0 & N/D & N/D \\
BMAD & 13 & 13 & 0 & N/D & N/D \\
Kiro & 39 & 39 & 0 & N/D & N/D \\
\bottomrule
\multicolumn{6}{p{\dimexpr\textwidth-2\tabcolsep}}{\footnotesize \textit{Nota:} Casos = pruebas automatizadas propias de cada variante (\textit{backend} y, cuando existen, \textit{frontend}). FAR = \textit{Functional Acceptance Rate}. N/D significa que no existe aún una medición común auditada, no que se hayan encontrado cero defectos.}
\end{tabularx}
\end{table}

\subsubsection{Dimensión E: Intervención Humana (Provisional)}

Los registros de invocaciones permiten contar solicitudes posteriores a la semilla inicial, pero no clasifican de manera uniforme cuáles fueron correcciones, pasos ordinarios del método o reintentos por fallas del servicio. Por ello la Tabla \ref{tab:intervencion_tetris_init} presenta este indicador observable sin llamarlo número definitivo de rondas de corrección. No se cronometró por separado el tiempo de revisión humana ni se administró un protocolo común para registrar defectos hallados exclusivamente por personas; ambos resultados son N/D. Entre las intervenciones documentadas constan la reparación del escenario imposible de Spec Kit, la reducción del control por teclado de Tessl a las tres flechas requeridas y la eliminación de metas cuantitativas no solicitadas en PRD, épicas y especificación de BMAD. Estos ejemplos ilustran el papel de la revisión de especificaciones, pero no constituyen un recuento completo de defectos.

\begin{table}[H]
\centering
\caption{Trazas de intervención observables -- Inicialización (Tetris Web)}
\label{tab:intervencion_tetris_init}
\begin{tabularx}{\textwidth}{l Y Y Y}
\toprule
\textbf{Herramienta} & \textbf{Solicitudes posteriores a la inicial} & \textbf{Tiempo de revisión humana} & \textbf{Defectos detectados solo por humanos} \\
\midrule
OpenSpec & 5 & N/D & N/D \\
Spec Kit & 5 & N/D & N/D \\
Tessl & 7 & N/D & N/D \\
BMAD & 8 & N/D & N/D \\
Kiro & 0 & N/D & N/D \\
\bottomrule
\end{tabularx}
\end{table}

